\documentclass[11pt]{article}

\usepackage{amsmath,amssymb,bm}
\usepackage{hyperref}
\usepackage{svg}
\usepackage[a4paper,margin=1in]{geometry}

\title{Two-Step OTFS I/O Equations}
\author{}
\date{}

\begin{document}
\maketitle

\section*{System model}

Consider a two-step OTFS modulator that uses an inverse symplectic finite
Fourier transform (ISFFT) followed by an OFDM modulator. Let \(N\) be the
number of subcarriers and \(M\) be the number of time slots, so that the
delay--Doppler grid contains \(NM\) information symbols
\(x[l,k]\), where \(l \in \{0,\dots,N-1\}\) is the delay index and
\(k \in \{0,\dots,M-1\}\) is the Doppler index. The delay and Doppler
resolutions on this grid are \(1/(N\Delta f)\) and \(1/(MT_S)\) respectively,
where \(\Delta f\) is the subcarrier spacing, \(T = 1/\Delta f\) is the useful
OFDM symbol duration, and \(T_S = T + T_{cp}\) is the OFDM symbol duration
including the cyclic prefix.

\subsection*{Transmitter: ISFFT \(+\) OFDM modulator}

The first step maps the delay--Doppler symbols \(x[l,k]\) to time--frequency
symbols \(X[n,m]\) on an \(N \times M\) grid using the ISFFT:
\begin{equation}
    X[n,m] = \frac{1}{\sqrt{NM}}
             \sum_{l=0}^{N-1}\sum_{k=0}^{M-1} x[l,k]\,
             e^{\,j2\pi\bigl(\tfrac{mk}{M} - \tfrac{nl}{N}\bigr)},
    \qquad
    \begin{aligned}
    &0 \leq n \leq N-1, \\
    &0 \leq m \leq M-1.
    \end{aligned}
\end{equation}
The symbols \(X[n,m]\) are then transmitted using the standard OFDM modulator,
i.e.\ \(X[n,m]\) is placed on the \(n\)th subcarrier of the \(m\)th OFDM
symbol. The \(m\)th OFDM symbol in the time domain is therefore
\begin{equation}
    s_m(t) = \sum_{n=0}^{N-1} X[n,m]\, e^{j2\pi n\Delta f\, t},
    \qquad 0 < t < T,
\end{equation}
and after appending a cyclic prefix the transmitted signal becomes
\begin{equation}
    s(t) = \sum_{m=0}^{M-1} s_m^{CP}(t - mT_S),
\end{equation}
exactly as in the OFDM case.

\subsection*{Channel}

The signal propagates through a \(P\)-path doubly selective channel with
delays \(\tau_p\), Doppler shifts \(\nu_p\), and complex gains \(\alpha_p\):
\begin{equation}
    r(t) = \sum_{p=0}^{P-1}\alpha_p\, s(t-\tau_p)\, e^{\,j2\pi \nu_p t} + z(t),
\end{equation}
with \(z(t)\) additive white Gaussian noise.

\subsection*{Receiver: OFDM demodulator \(+\) SFFT}

The receiver first performs OFDM demodulation. Following the derivation of
the OFDM I/O equation (and assuming the cyclic prefix is long enough to
absorb the channel delay spread so that inter--symbol interference is
avoided), the demodulator output on the \(n\)th subcarrier of the \(m\)th
OFDM symbol is
\begin{equation}
    Y[n,m] = \sum_{p=0}^{P-1} \tilde{\alpha}_p
             \sum_{n'=0}^{N-1} X[n',m]\,
             e^{-j2\pi n'\Delta f\, \tau_p}\,
             e^{\,j2\pi \nu_p\, mT_S}\,
             I_p(n'-n) + Z[n,m],
\label{eq:Y_general}
\end{equation}
where \(\tilde{\alpha}_p = \alpha_p\, e^{-j2\pi n\Delta f\, mT_S}\),
\(Z[n,m]\) is the projected noise, and
\begin{equation}
    I_p(n'-n) = \int_0^T e^{j2\pi(n'-n)\Delta f\, t'}\,
                e^{j2\pi \nu_p\, t'}\, dt'
              = \frac{e^{j2\pi \nu_p T} - 1}
                     {j2\pi\bigl[(n'-n)\Delta f + \nu_p\bigr]}
\end{equation}
captures the inter--carrier interference (ICI) caused by the Doppler shifts.
When the Doppler spread is small compared to \(\Delta f\) we have
\(I_p(n'-n) \approx T\,\delta[n'-n]\) and the simplified TF channel model
follows:
\begin{equation}
    Y[n,m] \;\approx\; \sum_{p=0}^{P-1} \tilde{\alpha}_p\, X[n,m]\,
                       e^{-j2\pi n\Delta f\, \tau_p}\,
                       e^{\,j2\pi \nu_p\, mT_S} + Z[n,m].
\label{eq:Y_simplified}
\end{equation}

The second step at the receiver is the SFFT, which maps the
time--frequency samples \(Y[n,m]\) back to the delay--Doppler domain:
\begin{equation}
    y[l',k'] = \frac{1}{\sqrt{NM}}
               \sum_{n=0}^{N-1}\sum_{m=0}^{M-1} Y[n,m]\,
               e^{-j2\pi\bigl(\tfrac{mk'}{M} - \tfrac{nl'}{N}\bigr)},
    \qquad
    \begin{aligned}
    &0 \leq l' \leq N-1, \\
    &0 \leq k' \leq M-1.
    \end{aligned}
\end{equation}

\section*{Delay--Doppler I/O equation (simplified, no ICI)}

Inserting~\eqref{eq:Y_simplified} into the SFFT and substituting the ISFFT
expression for \(X[n,m]\) gives
\begin{align}
y[l',k']
&= \frac{1}{\sqrt{NM}}\sum_{n=0}^{N-1}\sum_{m=0}^{M-1}
   \Biggl[
     \sum_{p=0}^{P-1} \tilde{\alpha}_p\, X[n,m]\,
     e^{-j2\pi n\Delta f\, \tau_p}\,
     e^{\,j2\pi \nu_p\, mT_S}
   \Biggr]
   e^{-j2\pi\bigl(\tfrac{mk'}{M} - \tfrac{nl'}{N}\bigr)} + \tilde{Z}[l',k']
   \notag\\[4pt]
&= \frac{1}{NM}\sum_{p=0}^{P-1}\tilde{\alpha}_p
   \sum_{l=0}^{N-1}\sum_{k=0}^{M-1} x[l,k]\,
   \underbrace{\sum_{n=0}^{N-1}
        e^{\,j2\pi n\bigl[\tfrac{l'-l}{N} - \Delta f\,\tau_p\bigr]}
   }_{\displaystyle A_p(l'-l)}\;
   \underbrace{\sum_{m=0}^{M-1}
        e^{\,j2\pi m\bigl[\tfrac{k-k'}{M} + T_S\,\nu_p\bigr]}
   }_{\displaystyle B_p(k-k')}
   + \tilde{Z}[l',k'],
\label{eq:OTFS_IO_general}
\end{align}
where \(\tilde{Z}[l',k']\) denotes the SFFT of the OFDM-domain noise.

\subsection*{Delay and Doppler kernels}

The two kernels are geometric sums:
\begin{align}
A_p(l'-l) &= \sum_{n=0}^{N-1}
             e^{\,j2\pi n\bigl[\tfrac{l'-l}{N} - \Delta f\,\tau_p\bigr]}
           = \frac{e^{\,j2\pi N\bigl[\tfrac{l'-l}{N} - \Delta f\,\tau_p\bigr]} - 1}
                  {e^{\,j2\pi \bigl[\tfrac{l'-l}{N} - \Delta f\,\tau_p\bigr]} - 1}, \\
B_p(k-k') &= \sum_{m=0}^{M-1}
             e^{\,j2\pi m\bigl[\tfrac{k-k'}{M} + T_S\,\nu_p\bigr]}
           = \frac{e^{\,j2\pi M\bigl[\tfrac{k-k'}{M} + T_S\,\nu_p\bigr]} - 1}
                  {e^{\,j2\pi \bigl[\tfrac{k-k'}{M} + T_S\,\nu_p\bigr]} - 1}.
\end{align}
They have the standard Dirichlet form and concentrate their energy at
\(l'-l \approx \Delta f\,\tau_p\, N\) and \(k'-k \approx T_S\,\nu_p\, M\),
i.e.\ at the normalized delay and Doppler indices of path \(p\).

\subsection*{Integer (on-grid) delay and Doppler}

In the special case where every path has its delay and Doppler aligned with
the delay--Doppler grid,
\begin{equation}
    \tau_p = \frac{l_p}{N\Delta f}, \qquad
    \nu_p  = \frac{k_p}{M T_S},
    \qquad l_p,\, k_p \in \mathbb{Z},
\end{equation}
the kernels collapse to discrete Kronecker deltas (with circular indexing):
\begin{align}
    A_p(l'-l) &= N\,\delta\!\left[(l - l' + l_p)\bmod N\right], \\
    B_p(k-k') &= M\,\delta\!\left[(k - k' + k_p)\bmod M\right].
\end{align}
Substituting these into \eqref{eq:OTFS_IO_general} eliminates the inner sums
and yields the well-known sparse delay--Doppler I/O relation:
\begin{equation}
    y[l',k'] = \sum_{p=0}^{P-1} \tilde{\alpha}_p\,
               x\!\left[(l' - l_p)_N,\, (k' - k_p)_M\right]
               + \tilde{Z}[l',k'],
\label{eq:OTFS_IO_integer}
\end{equation}
where \((\cdot)_N\) and \((\cdot)_M\) denote reduction modulo \(N\) and
\(M\). Each delay--Doppler tap therefore appears as a two--dimensional
circular shift of the input grid, scaled by the corresponding path gain.

\subsection*{Fractional delay and Doppler (general case)}

When \(\tau_p\) or \(\nu_p\) is not on the grid, the kernels do not collapse
to deltas and~\eqref{eq:OTFS_IO_general} retains its full form. Writing
\(\tau_p = (l_p + \kappa_p)/(N\Delta f)\) and
\(\nu_p  = (k_p + \iota_p)/(MT_S)\) with integer parts \(l_p,k_p\) and
fractional parts \(\kappa_p \in [-\tfrac12,\tfrac12)\),
\(\iota_p \in [-\tfrac12,\tfrac12)\), the kernels become
\begin{align}
A_p(l'-l) &= e^{\,j\pi(N-1)\bigl[\tfrac{l'-l-l_p}{N} - \tfrac{\kappa_p}{N}\bigr]}
             \cdot
             \frac{\sin\!\bigl(\pi(l'-l-l_p-\kappa_p)\bigr)}
                  {\sin\!\bigl(\tfrac{\pi}{N}(l'-l-l_p-\kappa_p)\bigr)}, \\
B_p(k-k') &= e^{\,j\pi(M-1)\bigl[\tfrac{k-k'+k_p}{M} + \tfrac{\iota_p}{M}\bigr]}
             \cdot
             \frac{\sin\!\bigl(\pi(k-k'+k_p+\iota_p)\bigr)}
                  {\sin\!\bigl(\tfrac{\pi}{M}(k-k'+k_p+\iota_p)\bigr)},
\end{align}
which are Dirichlet kernels centered at the integer delay/Doppler indices
\((l_p,k_p)\) and spread by the fractional parts \((\kappa_p,\iota_p)\).
Inserting these into~\eqref{eq:OTFS_IO_general} gives the full
delay--Doppler input--output relation including the inter--bin leakage that
is characteristic of fractional Doppler.

\section*{Compact form}

Combining the two kernels, define the effective delay--Doppler channel
response
\begin{equation}
    h_{\text{DD}}[l'-l,\,k'-k] = \frac{1}{NM}
        \sum_{p=0}^{P-1}\tilde{\alpha}_p\, A_p(l'-l)\, B_p(k-k').
\end{equation}
Then~\eqref{eq:OTFS_IO_general} can be written as a (twisted) 2D convolution
on the delay--Doppler grid:
\begin{equation}
    y[l',k'] = \sum_{l=0}^{N-1}\sum_{k=0}^{M-1}
               h_{\text{DD}}[l'-l,\,k'-k]\, x[l,k] + \tilde{Z}[l',k'].
\end{equation}
In the integer-tap case this reduces to~\eqref{eq:OTFS_IO_integer}, and the
sparsity of the delay--Doppler channel \(h_{\text{DD}}\) is what makes OTFS
attractive for high-mobility channels: each physical path produces a single
tap (up to leakage), whereas in the time--frequency domain it spreads over
the entire grid through the kernels \(I_p(\cdot)\).

\end{document}
